%% maths-tangente-calculatrice — écran de calculatrice graphique en mode « tangente ».
%% Reproduction neutre (aucune marque) de ce qu'affiche une calculatrice quand on lui
%% demande la tangente à une courbe : la courbe, la tangente tracée, le point de contact,
%% et une ligne de lecture qui donne x, f(x) et f'(x). La fonction est f(x) = x^2 + 2x + 3
%% et le point de contact a pour abscisse 2 : la machine affiche donc f(x) = 11 et
%% f'(x) = 6, de quoi retrouver y = 6x - 1 à la main.
%% L'équation de la tangente, elle, est enveloppée de \rappel : la variante muette
%% (« -muet ») ne montre que les trois nombres lus, c'est elle qui va dans les énoncés.
\documentclass[tikz,border=4pt]{standalone}
\usepackage{liaisons}
\begin{document}
\begin{tikzpicture}[
  boitier/.style={draw=black!55, line width=1pt, rounded corners=4pt, fill=black!2},
  barre/.style={fill=black!12, draw=black!35, line width=0.5pt},
  quad/.style={line width=0.35pt, draw=black!15},
  axeg/.style={line width=0.6pt, draw=black!55},
  lu/.style={font=\small}]

%% ---------------- le boîtier ----------------------------------------------
\draw[boitier] (0,0) rectangle (9,6);
\fill[barre] (0.1,5.2) rectangle (8.9,5.9);
\node[font=\small\bfseries, text=black!70, anchor=west] at (0.35,5.55) {GRAPHIQUE};
\node[font=\small, text=black!60, anchor=east] at (8.65,5.55) {tangente};
\draw[black!35, line width=0.5pt] (0.25,0.9) rectangle (8.75,5.05);

%% ---------------- la fenêtre graphique ------------------------------------
%% repère interne : 1 unité en x = 0,85 cm, 1 unité en y = 0,307 cm ;
%% l'origine du repère mathématique est placée en (5,35 ; 1,207).
\begin{scope}[shift={(5.35,1.207)}, x=0.85cm, y=0.307cm]
  \clip (-6,-1) rectangle (4,13);
  \foreach \xv in {-5,-4,-3,-2,-1,1,2,3} \draw[quad] (\xv,-1) -- (\xv,13);
  \foreach \yv in {2,4,6,8,10,12} \draw[quad] (-6,\yv) -- (4,\yv);
  \draw[axeg] (-6,0) -- (4,0);
  \draw[axeg] (0,-1) -- (0,13);
  \foreach \yv in {4,8,12} \node[font=\scriptsize, text=black!55, anchor=south west,
        inner sep=1pt] at (0.1,\yv) {$\yv$};
  \foreach \xv in {-4,2} \node[font=\scriptsize, text=black!55, anchor=north,
        inner sep=2pt] at (\xv,0) {$\xv$};
  %% la courbe et la tangente
  \draw[solideB, line width=1.4pt, line join=round]
       plot[domain=-4.32:2.32, samples=80, smooth] (\x,{\x*\x+2*\x+3});
  \draw[solideA, line width=1.2pt] (0,-1) -- (2.35,{6*2.35-1});
  \fill[solideA] (2,11) circle (2.6pt);
\end{scope}

%% ---------------- la ligne de lecture -------------------------------------
\fill[barre] (0.25,0.15) rectangle (8.75,0.75);
\node[lu, anchor=west] at (0.55,0.45) {$x=2$};
\node[lu, anchor=center] at (4.5,0.45) {$f(x)=11$};
\node[lu, anchor=east] at (8.45,0.45) {$f'(x)=6$};

%% ---------------- ce que la machine ne donne pas -------------------------
\rappel{\node[draw=black!45, fill=black!3, rounded corners=3pt, line width=0.5pt,
      inner sep=5pt, anchor=north, align=center, font=\small] at (4.5,-0.25)
     {équation de la tangente : $y=6x-1$};}
\end{tikzpicture}
\end{document}
